\chapter{Discussion and conclusion}
\label{chapter:discussion}

\section{Answering the scientific question}

\paragraph{Can PPO learn a low-thrust rendezvous guidance law?} Partially, and narrowly. On a
single, small, near-Earth synthetic target, PPO reliably learns a policy that rendezvous 20/20
times, beats both baselines, uses only 7.7\% more delta-v than a closed-form impulsive lower
bound, and -- critically -- remains successful when its exact thrust history is replayed through
an independent, higher-fidelity integrator (Section~\ref{sec:wp8}). That last result matters: it
rules out the possibility that stage 1's success is a numerical artifact of the RK4 propagator
used in training, and confirms the learned control law is solving the real physical problem.

But this success did not scale. Every attempt to move beyond the single easy target -- a larger
synthetic offset (stage 2), and real GTOC4 asteroids with a randomised target every episode (the
final multi-seed configuration) -- failed to beat a policy that simply coasts, across four
different targeted fixes for stage 2 and five independent seeds for the final configuration. The
failure-mode sweep in Section~\ref{sec:wp7} shows why this is not surprising in hindsight: the
stage-1 policy's success rate collapses from 10/10 to 0/10 the moment the target's semi-major-axis
offset is doubled, well short of stage 2's offset. So the honest answer is that PPO, with this
architecture, reward, and curriculum, learned to solve one problem instance rather than a
generalisable guidance law.

\paragraph{Does the learned controller resemble bang-bang optimal control?} In the narrow sense
of throttle, yes: the stage-1 policy thrusts at 100\% for every control step, with no coast phase
at all (Section~\ref{sec:wp7}). But the policy-probing result -- near-total insensitivity of the
commanded thrust direction to large, deliberate changes in the relative position error --
suggests this is not a generic closed-loop guidance law that happens to always want full thrust.
It looks more like a single, roughly time-indexed maneuver sequence memorised for one specific
transfer geometry, dressed up as a feedback policy by the fact that its inputs barely change from
the one situation it was trained on. Whether a policy trained with a broader curriculum (e.g.\
randomised $\Delta a$ even at stage-1 difficulty) would instead learn a genuinely reactive law is
an open question this project did not have time to test.

\section{Why generalisation failed: candidate mechanisms}

Several mechanisms, each independently documented in Chapter~\ref{chapter:results}, plausibly
contributed:

\begin{itemize}
  \item \textbf{Reward shaping's blind spot for loitering.} $\gamma=1$ potential-based shaping
        telescopes to net start-to-end progress; a policy that flies through the target's
        neighbourhood without settling gets almost the same shaped return as one that
        rendezvous. Stage 1 only converged after its tolerance was loosened enough that
        ``passing near enough'' became achievable; stage 2's tighter effective margin
        (delta-v margin only $1.2\times$) may have made the same failure mode unrecoverable
        rather than merely slow to fix.
  \item \textbf{Delta-v margin.} Both stage 2 and the real-asteroid pool were deliberately
        filtered/extended to keep at least a $1.5\times$ delta-v margin over the closed-form
        estimate, but that estimate is for an \emph{ideal} impulsive transfer; an imperfect
        learned policy that overshoots or corrects late has much less margin in practice.
  \item \textbf{Curriculum step size.} The jump from stage 1 ($\Delta a=0.02\,\mathrm{AU}$) to
        stage 2 ($\Delta a=0.1\,\mathrm{AU}$) is a $5\times$ increase, well past the point
        (Section~\ref{sec:wp7}) where the stage-1 policy's success rate had already collapsed to
        zero. A curriculum with intermediate steps was not tried, for lack of time.
\end{itemize}

\section{Limitations}

\begin{itemize}
  \item Stage 1's success is a single-seed result (Figure~\ref{fig:stage1-curve}); the final
        configuration's 5-seed multi-seed result is the more statistically meaningful comparison,
        and that one failed.
  \item The sensitivity sweep (Section~\ref{sec:wp6}) was run at half the training budget needed
        for the baseline itself to converge, so it mostly characterises difficulty-driven
        convergence speed under a fixed budget rather than sensitivity of converged performance.
  \item Stage 2 and stage 3/the final configuration remain open, unresolved failures: this report
        documents four and five falsified fix hypotheses respectively, but does not claim to have
        found the actual fix.
  \item The mass-optimality comparison (Section~\ref{sec:wp7}) uses a closed-form impulsive lower
        bound, not a true continuous-thrust optimal control solution; the $+7.7\%$ gap is
        indicative, not a tight bound.
\end{itemize}

\section{Future work}

The most direct next steps are: (1) re-run the sensitivity sweep at full training budget to
separate genuine sensitivity from convergence-speed artifacts; (2) a finer-grained curriculum
between stages 1 and 2, to test whether the generalisation cliff in
Figure~\ref{fig:failure-sweep} can be pushed outward gradually rather than jumped over; (3) a
non-telescoping or time-aware potential (e.g.\ penalising being far from the target at any point,
not just at episode end) to address the loitering blind spot directly; and (4) comparing the
learned trajectory against a true indirect (Pontryagin) or direct-transcription optimal-control
solution for the same stage-1 target, to turn the mass-optimality gap into a tight bound rather
than an approximate one.

\section{Conclusion}

A hand-written RK4 dynamics core, validated against Tudat to within a few metres once gravity
constants are matched, and a Gymnasium/PPO training pipeline, together learn a low-thrust
rendezvous policy that solves one specific target reliably, efficiently (within 7.7\% of an
impulsive delta-v lower bound), physically (verified under an independent integrator), and with a
fully bang-bang (always-on) thrust profile. That policy does not generalise: it fails sharply just
outside its training distribution, and every attempt to train directly on harder or randomised
targets -- across a curriculum stage and a five-seed final configuration -- failed to beat simply
coasting. The project's clearest scientific finding is therefore not that RL solves low-thrust
guidance, but a more specific and better-supported one: policy-gradient RL with potential-based
shaping can find an efficient, physically-verified solution to a \emph{single} low-thrust
rendezvous instance, without thereby learning a controller that generalises even modestly within
the same family of problems.
